Definition of supports 

``With FDM printers, geometric constraints relating to overhangs and base size can be partially addressed using automatically generated scaffolding known as supports and rafts''\cite[p.~391]{hudson_understanding_2016}

``many print centers enforced final check-overs by the operators on users’ designs and print settings before printing to reduce failed prints''\cite[p.~392]{hudson_understanding_2016}


\section{what does it solve}

\begin{itemize}
    \item what does the optimized workflow look like? 
    \item how it might make sharing files easier for beginners when supports and modifiers are included $\rightarrow$ lowering the access barrier for casuals, beginners and in educational contexts
    \item how it reduces cognitive workload $\rightarrow$ minimizing mistakes and waste due to prints going wrong 
\end{itemize}

``\emph{makers}, a community of enthusiasts who focus on fabrication, invention and experience sharing, and who collaborate and learn in environments known as \emph{makerspaces}''\cite{hudson_understanding_2016}

``thousands of public institutions, such as schools, libraries, and universities are now setting up print centers or creative hubs for end users to make use of digital fabrication technologies''\cite{hudson_understanding_2016}

``makerspaces that are often run by enthusiasts [\dots] attract engineers, entrepreneurs, inventors, hackers, craftsmen, and artists''\cite{hudson_understanding_2016}

``casual makers serve as a proxy for what it will be like for the general public to use 3D printing''\cite{hudson_understanding_2016}

``if HCI is going to be at the forefront of inventing new fabrication design tools and interfaces[\dots], we need to look beyong enthusiasts and understand the challenges and opportunities that exist in facilitating the interactions of casual makers with fabrication technology''\cite{hudson_understanding_2016}

``current tools that focus on supporting only one aspect of the workflow (\emph{\eg}, only modelling or only printing) may be less useful for casual makers and there is opportunity in HCI research to innovate on design tools that take into account the interdependencies within the 3D printing workflow''\cite{hudson_understanding_2016}

``while lowering the usability barriers of fabrication tools is important, it is perhaps even more important to weave-in expertise on best practices and troubleshooting tips within the walk-up-and-use tools used by casual makers to reduce theiir dependecy on operators''\cite{hudson_understanding_2016}

``opportunities that exist in building the next generation of inter-connected fabrication tools with features for supporting expertise sharing''\cite{hudson_understanding_2016}

``enthusiasts and professional users—that is, they were trained or technically literate in fabrication''\cite{hudson_understanding_2016}


% 3mf 

    \begin{table}
        \renewcommand{\arraystretch}{1.3}
        \begin{tabular}{m{0.21\linewidth} m{0.5\linewidth} m{0.14\linewidth}}
            \textbf{Name} & \textbf{Description} & \textbf{Required /} \textbf{Optional} \\
            \hline
            3D Model & Contains the description of one or more 3D objects for
manufacturing. & Required \\
            Core Properties & The OPC part that contains various document properties. & Optional \\
            Object Thumbnail & Contains a small JPEG or PNG image that represents a 3D object in a 3D Model. & Optional \\
            3D Texture & Contains a texture used to apply color to a 3D object in the 3D Model part (available for extensions) & Optional \\
            PrintTicket & Provides settings to be used when outputting the 3D object(s) in the 3D Model part. & Optional \\
            Custom Parts & OPC parts that are associated with metadata & Optional
        \end{tabular}\smallskip

        \caption{Excerpt of 3MF document parts\cite[p.~4]{3mf_consortium_3mf_2025}}
        \label{doc_parts}
    \end{table}

    \section{Fusion Add-In}

\begin{itemize}
    \item what is CAD software 
    \item a lil about their API: add-in vs. script, Python vs. C++
    \item why I chose to realise it as a Python add-in
    \item possibly different approaches if more than one way to solve it came up 
\end{itemize}

As mentioned before, the software supports developers to create additional functionalities by providing an extensive API and the option to create either an add-in or a script directly from within Fusion's UI. 
Using this, Fusion sets up the project structure as well as some example code to make it easier to get started especially for beginners. 
Fusion runs natively on Python but offers to create the add-in or script either in Python or in C++. 

\subsection{Script vs. Add-in}\cite{autodesk_inc_fusion_Script_addin_2025}
\begin{itemize}
    \item \textbf{Run on Startup} – This setting is add-in-specific and indicates if the add-in should be run automatically when Fusion is started.  Most add-ins will want to take advantage of this capability so the commands they define will be available to the user as soon as Fusion starts.
    \item The user executes a script through the “Scripts and Add-Ins” command, and the execution stops immediately after the run function is complete. A script runs, and then it is done.
    \item An add-in is typically automatically loaded by Fusion when Fusion starts up. An add-in also usually creates one or more custom commands and adds them to the user interface during startup. The add-in continues to run throughout the Fusion session, so it can react whenever the user executes any of its commands. The add-in remains running until Fusion is shut down or the user explicitly stops it through the “Scripts and Add-Ins” dialog. When it stops, it cleans up whatever user interface customization it created in its stop function.
\end{itemize}

\subsection{Python vs. C++}



\chapter{Problem}\label{ch:problem}

As shown in the previous chapters \ref{chap:apply-materials} and \ref{chap:apply-settings}, there are a lot of parameters to be repeatedly applied before a 3D model can be printed. 
So, why would this be a problem? 
There are in fact multiple points to consider regarding the necessity of repetitive manual steps in a common digital workflow especially in educational and casual makers' contexts. 

``Current computer-aided design and manufacturing (CAD-CAM) software which assume a well-understood fabrication process are mismatched to a material-driven trial-and-error process. This creates friction when exploring—copying files, loading settings—but also produces downstream effects on how designs are stored, shared, and replicated.Slight variation in material composition and machine state affect physical outcomes''[its not the shape]

``We identify two important limitations of current fabrication software in material-driven practices: the ability to quickly explore a material’s design space, and the ability to document the insights gained through this exploration.''[its not the shape]

``tuning settings allows practitioners to experiment with surface texture and aesthetic properties, control material properties like stiffness, and optimize print settings''[same]

\section{Negative Effects of Cognitive Workload}

Setting the parameters once might be an engaging part of the entire process of getting to the final product. 
Especially in an educational context, the users do learn a lot about the possibilities the slicer and printer offer, the entire fabrication process and how to approach a project on their own. 
Doing the exact same steps repeatedly however, does not teach the users anything, while still requiring a certain amount of cognitive workload. 

What does this even mean? 
As Kosch \etal summarized in \cite{kosch_survey_2023}, there are various definitions used in the field of human computer interaction (HCI).
One of them differentiates between \textit{intrinsic} and \textit{extraneous load}.
The former encompasses the ``inherent complexity of a task'' and therefore ``depends on the users individual cognitive resources and proficiency'', i.e. their skills. 
In other words, how difficult and extensive is a task to begin with and how well equipped is the person in dealing with the complexity. 
On the other hand, the latter describes cognitive resources needed due to the ``task representation that can be perceived'' by the users.
For instance, how does the design of a user interface either reduce or increase friction within a given workflow?

Cognitive workload being a part of the workflow isn't a negative aspect in itself. 
To a certain degree it is even necessary to foster task engagement which increases the learning effects and the feel of accomplishment.
In the process at hand both types of cognitive workload can be identified.
When initially setting the printing parameters, the task will be inherently complex; more or less so, depending on the complexity of the project itself. 
Each repetition however does turn into extraneous load, since having done the steps once before, the task requires focus without any added value of engagement, thus creating additional friction by way of design. 
One might argue that for beginners as well as casual users the step can nevertheless be overwhelming within the larger context of several new and complex tasks while for advanced makers it is simply tedious and time consuming. 

Due to the task being only one part of a multifaceted list of steps, this kind of friction claims resources that won't be free to solve other parts like setting up the printer. 
One might argue that people should just focus on one task at a time, plan ahead and not try to multitask, however with the added time constraint oftentimes present in particular in educational environments but also when using community equipment (e.g. in makerspaces), it becomes a necessary evil. %% confusing additional argument - either expand on it or throw it out 

%% BIG FLAW: where's the evidence / metrics to support the thesis of this? Can I say "these statements are largely based on analytical observations of classroom scenarios since the scope of the project did not extent to user studies"?

$\rightarrow$ increasing amount of errors? 

``In general, usability is an essential HCI concept and is concerned with making systems easy to learn, easy to use, and with minimal error frequency and severity.'' [Sustainable Design]

TBC

\section{Sharing Models Online}

In their study of newcomers to 3D printing, Hudson \etal worked out a few key problems impacting the whole user experience, one of them being the problem of settings causing print failures.\cite{hudson_understanding_2016} 
While this can't be entirely avoided through the solution presented in the following chapter, it can be reduced. 
Since creating a 3D model does present an obstacle, people turn ``to premade 3D models available through services such as \textit{Thingiverse.com}''.\cite{hudson_understanding_2016}

However even when using such pre-made models, which oftentimes come with elaborate instructions regarding the printing parameters, users were successful only with the help of print center operators. 
Especially when 3D printers are provided by community institutions like libraries though, the staff might not be trained in the use of these printers as well as in makerspaces.
This results in failed prints, frustrating the user but also creating material waste. 


``With personal fabrication, however, the responsibility of waste generation is shifting directly to individual makers. As making becomes more accessible, the amount of material waste generated at the grassroots level is increasing.'' [Make Making Sustainable]

``the responsibility of waste generation is shifting directly to individual makers''[dito]

``Almost all the makers discussed how their projects require design iterations and prototyping. These processes naturally generate a number of intermediate prototypes, which, as explained in Section 4.1.2, serve as a natural source of waste. Thus, optimizing design iterations and using fewer prototypes, ideally without compromising design goals, will result in reduced material usage, fewer physical artifacts, and more sustainable actions.''[dito]

TBC


% Negative Effects resulting from manual Repetitive Steps 
%   increased amount of errors due to necessary repetitive steps
%   higher frustration levels
%   discouraged beginners, waste of time / material / resources in early stages of teaching 
%   ``If the user’s print failed or did not match their design intent, they would have to adjust either their model or print settings to correct the issue, often with the operator’s assistance''\cite{hudson_understanding_2016} $\rightarrow$ frustration tolerance 

%. the need for elaborate and extensive printing instructions when sharing a complex model online 

``personal fabrication introduces a distinct challenge that was not present during the democratization of personal computing: the environmental impact of material waste produced by makers''[make making sustainable]

\section{Fabrication Documentation}

``Documenting fabrication workflows is a diffcult and time-consuming activity. Practitioners often have to choose between making progress on their project and creating documentation''[its not the shape]

``most readers of documentation on the website Instructables are interested in customization, not rote recreation''[same]

``Currently when sharing 3D designs online, settings are optionally and manually disclosed alongside models.''[same]